\subsection{PROPERTIES OF EQUALITY}

From now on we shall consider only equalitarian theories. Let \(\mathfrak{G}\) be such a theory, and let \(\mathfrak{G}_0\) be the theory whose signs are those of \(\mathfrak{G}\) and whose only axioms are those provided by schemes S1 through S7. The theory \(\mathfrak{G}_0\) is weaker than \(\mathfrak{G}\) (§2, no. 4) and has no constants. The following three theorems are theorems in \(\mathfrak{G}_0\) .

THEOREM 1. x = x.

Let \(S\) denote the relation \(x = x\) in \(\mathcal{C}_0\) . By C27 (§4, no. 1), for every relation \(R\) in \(\mathcal{C}_0\) , \((\forall x)(R \leftrightarrow R)\) is a theorem in \(\mathcal{C}_0\) , and therefore, by S7, \(\tau_x(R) = \tau_x(R)\) , that is to say \((\tau_x(R)|x)S\) , is a theorem in \(\mathcal{C}_0\) . Taking \(R\) to be the relation ``not \(S\) '' and considering C26 (§4, no. 1), we see that \((\forall x)S\) is a theorem in \(\mathcal{C}_0\) . By C30 (§4, no. 3), \(S\) is therefore a theorem in \(\mathcal{C}_0\) .

The relation \((\forall x)(x = x)\) is also a theorem in \(\mathfrak{C}_0\) ; and if \(T\) is a term in \(\mathfrak{C}_0\) , then \(T = T\) is a theorem in \(\mathfrak{C}_0\) (cf.~§4, no. 3). It is possible to transform later theorems in the same way into theorems in which no letter appears or into metamathematical criteria. From now on we shall not explicitly perform these transformations, but we shall often implicitly make use of them.

THEOREM 2. \((x = y) \Longleftrightarrow (y = x)\) .

Suppose that the relation \(x = y\) is true. By S6, the relation

\[
(x = y) \Longrightarrow ((x | y) (y = x) \Longleftrightarrow (y | y) (y = x)),
\]

that is

\[
(x = y) \Longrightarrow ((x = x) \Longleftrightarrow (y = x)),
\]

is true. Therefore \((x = x) \Longleftrightarrow (y = x)\) is true. By Theorem 1 it follows that y = x is true, and the theorem is proved.

THEOREM 3. \(((x = y)\) and \((y = z))\Longrightarrow (x = z)\) .

Let us adjoin the hypotheses \(x = y\) , \(y = z\) to the axioms of \(\mathfrak{G}_0\) . By S6 the relation \((x = y) \Longrightarrow ((x = z) \Longleftrightarrow (y = z))\) is true. Hence

\[
(x = z) \Longleftrightarrow (y = z),
\]

and consequently \(x = z\) , are true.

C44. Let \(\pmb{x}\) be a letter and let \(\pmb{T}, \pmb{U}, V\{\pmb{x}\}\) be terms in \(\mathfrak{C}_0\) . Then the relation \((\pmb{T} = \pmb{U}) \Longrightarrow (V\{\pmb{T}\} = V\{\pmb{U}\})\) is a theorem in \(\mathfrak{C}_0\) .

For let y and z be two distinct letters which are distinct from x and from the letters which appear in T, U, V. Adjoin the hypothesis y = z. Then, by S6,

\[
((y | z) (V \{y \} = V \{z \})) \Longleftrightarrow (V \{y \} = V \{z \}),
\]

that is to say \((V\{y\} = V\{y\}) \Longleftrightarrow (V\{y\} = V\{z\})\) , is true. Now, \(V\{y\} = V\{y\}\) is true by Theorem 1; hence \(V\{y\} = V\{z\}\) is true.

From all this it follows that \((y = z) \Longrightarrow (V\{y\} = V\{z\})\) is a theorem in \(\mathfrak{G}_0\) , say \(A\) . But \((T|y)(U|z)A\) is precisely

\[
(T = U) \Longrightarrow (V \{T \} = V \{U \}).
\]

A relation of the form \(T = U\) , where \(T\) and \(U\) are terms in \(\mathfrak{G}\) , is called an equation; a solution (in \(\mathfrak{G}\) ) of the relation \(T = U\) , considered as an equation in a letter \(x\) , is therefore (§2, no. 2) a term \(V\) in \(\mathfrak{G}\) such that \(T\{V\} = U\{V\}\) is a theorem in \(\mathfrak{G}\) .

Let \(\pmb{T}\) and \(\pmb{U}\) be two terms in \(\mathfrak{C}\) , and let \(x_1, x_2, \ldots, x_n\) be the letters which appear in \(\pmb{T}\) but not in \(\pmb{U}\) . If the relation

\[
(\exists x _ {1}) \dots (\exists x _ {n}) (T = U)
\]

is a theorem in \(\mathfrak{C}\) , we say that \(U\) can be put in the form \(T\) (in \(\mathfrak{C}\) ). Let \(R\) be a relation in \(\mathfrak{C}\) and let \(y\) be a letter. Let \(V\) be a solution (in \(\mathfrak{C}\) ) of \(R\) , considered as a relation in \(y\) . If every solution (in \(\mathfrak{C}\) ) of \(R\) , considered as a relation in \(y\) , can be put in the form \(V\) , then \(V\) is said to be the complete solution (or general solution) of \(R\) (in \(\mathfrak{C}\) ).

